\begin{problem}[第32届全国中学生物理竞赛决赛]{52}
如图，一个上端固定、内半径为 $R_{1}$ 的玻璃圆筒底部浸没在大容器中的水面之下，未与容器底接触；将另一个外半径为 $R_{2}$（$R_{2}$ 略小于 $R_{1}$）的同质玻璃制成的圆筒（上端固定）放置其中，不接触容器底，保持两玻璃圆筒的中轴线重合且竖直。设水与玻璃之间的接触角为 $\theta$（即气-液界面在水的最高处的切线与固-液界面的夹角，见局部放大图）。已知水的密度和表面张力系数分别为 $\rho$ 和 $\alpha$，地面重力加速度的大小为g。
\begin{subproblem}
在毛细现象中，系统毛细作用引起的势能变化，可以等效地看作固-液-气分界线上液体表面张力在液面上升或下降过程中做的功。试证：以大容器中的水面处为系统势能零点，则当水由于毛细作用上升高度为 h 时，系统毛细作用的势能为
\begin{equation}
E_{\mathrm{s}}(h) = -2\pi\alpha\cos\theta\left(R_{1}+R_{2}\right)h;
\end{equation}
\end{subproblem}
\begin{subproblem}
试导出系统在水由于毛细作用上升的过程中释放的热量；
\end{subproblem}
\begin{subproblem}
如果在超万米高空飞机中短时间内做此实验，飞机可视为做匀速直线运动。测得该系统在同样过程中放出的热量为 $Q$，试估计此时飞机距离地面高度。假定该飞机内水的表面张力系数和接触角与地面情形相同。

以上计算不考虑地球自转。
\end{subproblem}
\end{problem}